\documentclass[12pt, a4paper]{article}

% --- PAQUETES NECESARIOS ---
\usepackage[utf8]{inputenc}      % Permite usar caracteres en español como tildes y ñ
\usepackage[spanish]{babel}      % Configura el idioma español para títulos automáticos y guiones
\usepackage{geometry}            % Para ajustar los márgenes
\usepackage{amsmath}             % Para fórmulas matemáticas
\usepackage{parskip}             % Usa espacio vertical entre párrafos en lugar de sangría
\usepackage{graphicx}            % Para incluir imágenes (logo)
\usepackage{hyperref}            % Para crear hipervínculos (útil en PDF)

% --- CONFIGURACIÓN DE MÁRGENES ---
\geometry{
    a4paper,
    left=2.5cm,
    right=2.5cm,
    top=2.5cm,
    bottom=3cm
}

% --- METADATOS DEL DOCUMENTO ---
\title{
    \Huge Guion para Video de Filosofía \\
    \vspace{0.5cm}
    \Large \textbf{"Piketty y la Máquina de la Desigualdad"}
}
\author{Basado en la rúbrica del Departamento de Filosofía \\ Instituto Nacional - 2025}
\date{\today}

% --- INICIO DEL DOCUMENTO ---
\begin{document}

\maketitle
\thispagestyle{empty} % No numera la página del título
\tableofcontents
\newpage

% --------------------------------------------------------------------
\section{Introducción: El Gancho (00:00 - 00:30)}
% --------------------------------------------------------------------

\textbf{Finalidad:} Captar la atención del espectador y presentar el tema de forma directa. Esto apunta al criterio \textit{Creatividad y comunicación audiovisual}.

\vspace{0.5cm}

\textbf{Voz en Off (enérgica y clara):}
"¿Alguna vez te has preguntado por qué, en un mundo con tanta riqueza, la balanza parece siempre inclinarse hacia los mismos? ¿Es justo que el lugar donde naces determine tus oportunidades en la vida? Hoy no hablaremos de suerte, hablaremos de un sistema."

\vspace{0.5cm}

\textit{\textbf{Sugerencia Visual:}} 
\begin{itemize}
    \item Imágenes de alto contraste: un rascacielos de lujo junto a una vivienda precaria.
    \item Gráficos animados que muestran cómo un pequeño porcentaje de la población acumula la mayor parte de la riqueza mundial.
    \item Música de fondo intrigante y moderna.
    \item Un miembro del equipo con voz clara y buena dicción graba esta introducción para darle fuerza.
\end{itemize}

% --------------------------------------------------------------------
\section{¿Quién es Thomas Piketty? (00:30 - 01:15)}
% --------------------------------------------------------------------

\textbf{Finalidad:} Presentar al autor de forma concisa, demostrando una comprensión inicial de su perfil.

\vspace{0.5cm}

\textbf{Narrador 1:}
"Para entender este sistema, debemos conocer a uno de sus críticos más influyentes: Thomas Piketty, un economista francés nacido en 1971. Con un doctorado de la London School of Economics, Piketty no es solo un académico; es un detective de datos."

\vspace{0.5cm}

\textbf{Narrador 2:}
"Ha dedicado su vida a analizar siglos de información económica para responder una pregunta fundamental: ¿De dónde viene la desigualdad y por qué persiste?"

\vspace{0.5cm}

\textit{\textbf{Sugerencia Visual:}} 
\begin{itemize}
    \item Aparece una foto de Thomas Piketty. A su lado, una línea de tiempo animada con sus hitos: École Normale Supérieure, LSE, EHESS.
    \item Portadas de sus libros más importantes ("El capital en el siglo XXI", "Capital e ideología") aparecen en pantalla.
\end{itemize}


% --------------------------------------------------------------------
\section{La Fórmula de la Desigualdad: \( r > g \)}
% --------------------------------------------------------------------
\textit{Duración estimada: 01:15 - 02:45}

\textbf{Finalidad:} Explicar con claridad la idea central del filósofo usando ejemplos (clave para \textit{Comprensión del pensamiento filosófico}).

\vspace{0.5cm}

\textbf{Narrador 1:}
"En su libro más famoso, 'El capital en el siglo XXI', Piketty llega a una conclusión impactante, que se puede resumir en una simple fórmula: \( r > g \)."

\vspace{0.5cm}

\textbf{Narrador 2:}
"Esto significa que la tasa de rendimiento del capital (la 'r'), es decir, el dinero que se gana invirtiendo, históricamente crece más rápido que la economía en general (la 'g'), que representa el crecimiento de los salarios y la producción. En palabras simples: es más fácil hacer dinero cuando ya tienes dinero, que trabajando."

\vspace{0.5cm}

\textbf{Narrador 1:}
"Pensemos en un ejemplo actual: una persona hereda un gran portafolio de inversiones que le rinde un 5\% anual. Otra persona trabaja y recibe un aumento de sueldo del 2\% anual, que es el crecimiento promedio de la economía. Con el tiempo, la brecha entre ellos no se cierra, sino que se hace cada vez más grande. Esto, según Piketty, concentra la riqueza y el poder en pocas manos, creando un riesgo de plutocracia: el gobierno de los ricos."

\vspace{0.5cm}

\textit{\textbf{Sugerencia Visual:}} 
\begin{itemize}
    \item La fórmula \(r > g\) aparece en grande. "r" se ilumina y muestra íconos de inversiones, acciones, propiedades. "g" se ilumina y muestra íconos de salarios y producción.
    \item Una animación simple de dos personas: una (inversionista) ve su montaña de monedas crecer rápidamente. La otra (trabajador) ve su montaña crecer mucho más lento.
\end{itemize}


% --------------------------------------------------------------------
\section{La Justificación: ¿Por qué aceptamos la desigualdad?}
% --------------------------------------------------------------------
\textit{Duración estimada: 02:45 - 04:00}

\textbf{Finalidad:} Profundizar en el análisis, abordando las ideologías que sustentan el sistema. Cumple con "Muestra comprensión profunda" de la rúbrica.

\vspace{0.5cm}

\textbf{Narrador 2:}
"Pero si esto es tan evidente, ¿por qué lo aceptamos? Piketty argumenta en 'Capital e ideología' que toda sociedad crea una historia para justificar sus desigualdades. Nuestra historia actual se basa en una idea muy poderosa: la \textbf{meritocracia}."

\vspace{0.5cm}

\textbf{Narrador 1:}
"Se nos dice que vivimos en una sociedad justa, donde todos tenemos las mismas oportunidades y el éxito es resultado del esfuerzo y el talento. Si eres rico, es porque lo mereces. Si eres pobre... es tu culpa. Pero, ¿es esto realmente cierto?"

\vspace{0.5cm}

\textit{\textbf{Sugerencia Visual:}} 
\begin{itemize}
    \item La palabra "MERITOCRACIA" aparece en pantalla, como si fuera el título de un libro sagrado.
    \item Imágenes de graduaciones, ascensos laborales, gente celebrando el éxito.
\end{itemize}

% --------------------------------------------------------------------
\section{Análisis y Reflexión Personal: Rompiendo el Mito}
% --------------------------------------------------------------------
\textit{Duración estimada: 04:00 - 05:30}

\textbf{Finalidad:} Conectar las ideas de Piketty con la realidad propia, cumpliendo el requisito de \textit{Análisis y reflexión personal} para la máxima puntuación.

\vspace{0.5cm}

\textbf{Narrador 1:}
"Aquí es donde el pensamiento de Piketty nos interpela directamente. ¿Hasta qué punto la meritocracia es justa en nuestra realidad? Pensemos en el acceso a la educación. En teoría, todos pueden acceder a una buena universidad, pero en la práctica, ¿un estudiante de un colegio con menos recursos compite en igualdad de condiciones con alguien que tuvo acceso a preuniversitarios, clases particulares e idiomas desde niño?"

\vspace{0.5cm}

\textbf{Voz de un miembro del equipo 1 (tono reflexivo):}
"Piketty nos hace cuestionar si el 'mérito' no depende, en gran parte, de las herramientas que te entregan al nacer. La desigualdad de la que habla no es solo económica; es una desigualdad de oportunidades que el sistema tiende a presentar como justa."

\vspace{0.5cm}

\textbf{Voz de un miembro del equipo 2 (tono analítico):}
"Relacionando esto con nuestra experiencia, vemos cómo las redes de contacto o el capital cultural de una familia influyen mucho más que el simple esfuerzo. La idea de Piketty no es eliminar el mérito, sino preguntarnos cómo crear un punto de partida que sea verdaderamente justo para todos."

\vspace{0.5cm}

\textit{\textbf{Sugerencia Visual:}} 
\begin{itemize}
    \item Imágenes de diferentes realidades educativas en Chile.
    \item Un mapa de Chile que muestre la desigualdad de acceso a la educación superior por región o nivel socioeconómico.
    \item Los miembros del equipo pueden grabarse a sí mismos (o solo sus voces) para esta sección.
\end{itemize}

% --------------------------------------------------------------------
\section{Conclusión (05:30 - 06:00)}
% --------------------------------------------------------------------

\textbf{Finalidad:} Cerrar el video con una síntesis clara y una pregunta que invite a la reflexión.

\vspace{0.5cm}

\textbf{Narrador 2:}
"En resumen, Thomas Piketty nos muestra con datos que la desigualdad no es un fenómeno natural, sino el resultado de estructuras políticas e ideológicas. Su fórmula \( r > g \) expone la mecánica del sistema, y su crítica a la meritocracia nos obliga a cuestionar las historias que nos contamos sobre el éxito y el fracaso."

\vspace{0.5cm}

\textbf{Narrador 1:}
"La pregunta final no es si la desigualdad existe, sino qué tipo de sociedad queremos construir. Una que la justifique, o una que se atreva a imaginar un futuro más equitativo. Y tú, ¿qué piensas?"

\vspace{0.5cm}

\textit{\textbf{Sugerencia Visual:}} 
\begin{itemize}
    \item Transición rápida de todos los conceptos clave del video.
    \item La pantalla se va a negro y queda la pregunta final: "¿Qué sociedad queremos construir?".
    \item Créditos del equipo y música final.
\end{itemize}

\end{document}
% --- FIN DEL DOCUMENTO ---
